\documentclass[a4paper]{scrartcl}
\usepackage{fancyhdr}
\pagestyle{fancy}
\usepackage[english]{babel}
\usepackage[utf8]{inputenc}
\usepackage{graphicx}
\usepackage{url}
\usepackage{textcomp}
\usepackage{amsmath}
\usepackage{lastpage}
\usepackage{pgf}
\usepackage{wrapfig}
\usepackage{fancyvrb}

% Listing header

\usepackage{listings}
\usepackage{xcolor}

\definecolor{lightgray}{gray}{0.95}

\lstdefinestyle{javastyle}{
    backgroundcolor=\color{lightgray},
    basicstyle=\ttfamily\small,
    numbers=left,
    numberstyle=\tiny,
    numbersep=8pt,
    frame=single,
    showstringspaces=false,
    breaklines=true,
    breakautoindent=true,
    breakindent=2em,
    keepspaces=true,
    columns=fullflexible,
    tabsize=4,
    language=Java,
    keywordstyle=\color{blue},
    commentstyle=\color{gray!70},
    stringstyle=\color{orange!90!black}
}

\lstdefinestyle{outputstyle}{
    basicstyle=\ttfamily\small,
    frame=single,
    breaklines=true,
    columns=fullflexible,
    keepspaces=true,
    showstringspaces=false
}

% Create header and footer
\headheight 27pt
\pagestyle{fancyplain}
\lhead{\footnotesize{Object-Oriented Design, IV1350}}
\chead{\footnotesize{Seminar 2 Solution}}
\rhead{}
\lfoot{}
\cfoot{\thepage\ (\pageref{LastPage})}
\rfoot{}

% Create title page
\title{Seminar 3}
\subtitle{Object-Oriented Design, IV1350}
\author{Mo Wang mowang@ug.kth.se}
%\date{\today} Prints today's date
\date{2026-04-01}

\begin{document}

\maketitle
\noindent\textbf{Project Members:} \\ \hfill
% \tableofcontents %Generates the TOC
[Mo Wang, mowang@ug.kth.se]

\section*{Declaration:}

By submitting this assignment, it is hereby declared that all group members listed above have contributed to the solution. It is also declared that all project members fully understand all parts of the final solution and can explain it upon request.

It is furthermore declared that no part of the solution has been copied from any other source (except for resources presented in the Canvas page for the course IV1350), and that no part of the solution has been provided by someone not listed as a project member above.

\section{Introduction}
\label{sec:intro}

This report delves into the implementation and testing of the \textit{Repair Electric Bike} system. The purpose of Seminar~3 is to translate the design produced in Seminar~2 into executable source code and to verify the program behavior using unit tests. The implementation process follows the guidelines of object-oriented programming and testing, which are described in Chapters~6 and~7 of \textit{A First Course in Object-Oriented Development}.

The implementation follows the previously designed class diagrams and communication diagrams as an initial structure. During implementation, additional decisions regarding encapsulation, cohesion, coupling, naming, and code structure are continuously made in order to produce readable, maintainable, and robust Java code.

Automated unit testing is carried out in parallel with the implementation that goes parallel to the implementation. Tests are implemented using JUnit to validate the behavior of the core parts of the system and to ensure that all parts of the program work as intended when changes are made. The tests are designed according to the testing principles presented in Chapter~7 of the course book, emphasizing self-evaluating tests and meaningful coverage of program behavior.

\subsection{Requirements of Seminar 3}

Seminar~3 consists of two main parts: program implementation and unit testing. Unit tests are carried out incrementally after implementing each system operation.

The first part is to implement the basic flow and the startup sequence of the Repair Electric Bike system that was designed in Seminar~2. The implementation follows these requirements:

\begin{itemize}
    \item The program is written in Java and is compilable and executable.
    \item Only the basic workflow and startup sequence are implemented while alternative flows are intentionally omitted in this seminar.
    \item The user interface is replaced by a single \texttt{View} class containing hard-coded calls to the controller, which prints all returned results to \texttt{System.out}, including the repair order printout.
    \item No database or external systems are implemented. Instead, the integration layer stores data in memory using placeholder objects.
    \item Exceptions are not used, since exception handling is introduced in a later seminar.
    \item The code follows the programming guidelines described in Chapter~6 of the course book, including naming conventions, high cohesion, low coupling, good encapsulation, and one Javadoc comment for each public declaration.
\end{itemize}

The second part is to write unit tests. The testing requirements are as follows:

\begin{itemize}
    \item Unit tests are implemented using the JUnit framework.
    \item Tests cover all relevant classes in the controller, model, and integration
          layers.
    \item Simple DTO classes, as well as methods consisting only of constructors,
          setters, and getters, are excluded from testing.
    \item Methods that produce output to \texttt{System.out} are not tested, in
          accordance with the seminar requirements.
\end{itemize}

\subsection{Implementing Code Based on Design Diagrams}

The class diagrams and system operation diagrams from Seminar~2 build the foundation for implementing the program structure and behavior. The class diagram provides a blueprint for the overall architecture and package structure, while the communication diagrams provide guidance for implementing of system operations by explicitly describing message flow and object interaction.

By following these diagrams during implementation, the intended execution flow of each system operation is preserved and ambiguity is reduced. At the same time, the implementation is continuously evaluated against object-oriented design principles to ensure that the resulting code remains cohesive, encapsulated, and loosely coupled.

\section{Method}

% \textbf{This chapter tells \textit{how} you solved the task.} Explain how you worked and what tools you used when solving the tasks. \textit{Do not explain your solution and do not present entire sub-results in detail. It's enough to explain which steps you performed and perhaps give a short example.} For example, some of the things to explain in the report for seminar one, is that you used a category list and noun identification to find class candidates for the domain model, and to tell which UML editor you used.
% Goal: Explain how you worked (not the solution itself)

% Slight help:
% In the Method chapter of this report, the development process must be explained, including how code structure, layering, and design decisions were made. 
% The Method chapter must explain how the tests were designed and how test cases were selected.

This chapter describes the development process for implementing and testing the \textit{Repair Electric Bike} system in Seminar~3. It explains how the design diagrams were translated into executable Java code, how unit tests were carried out, and which tools and guidelines were used.

\subsection{Implementation Approach}

The implementation process started by translating the design diagrams produced in Seminar~2 into an initial Java code structure. The class diagram was used to define the overall architecture, including package structure and class responsibilities, while the communication diagrams guided the incremental implementation of each system operation.

Development began by creating package structures and class skeletons without functionality, according to the class diagrams. These skeletons were then extended iteratively, implementing one system operation at a time. Special attention was given to the startup sequence to ensure that all necessary objects were created and connected correctly before executing system operations. After each system operation, tests were written to detect potential errors, which were resolved before continuing with the next system operation.

Throughout the implementation, object-oriented programming guidelines from Chapter~6 of
the course book were continuously applied. Design decisions regarding encapsulation,
cohesion, and coupling were revisited when needed, and minor design adjustments were made
whenever implementation revealed potential improvements.

\subsection{Dividing Code into Packages}

The code was organized into packages corresponding to the layered MVC architecture defined
during design. Separate packages were created for the controller, model, integration, and
startup layers, each with a clearly defined responsibility.

Classes that did not naturally belong to a specific layer, such as utility functionality
used across multiple layers, were placed in a separate \texttt{util} package. This
organization improved readability and helped maintain low coupling between layers while
preserving high cohesion within each package.

The separation into layers produced low coupling and high cohesion, as well as improved readability, maintainability, and extensibility, since each package has a well-defined responsibility.

\subsection{Following code conventions}

During implementation, Java code conventions were followed consistently to improve readability and maintainability. Package names are written in lowercase and structured to reflect project organization, while class and interface names use descriptive identifiers written in \texttt{PascalCase}.

Method, field, and variable names follow \texttt{camelCase} naming conventions and are chosen to clearly describe their purpose. Constants declared as \texttt{static final} fields are written in uppercase letters with words separated by underscores. Adhering to these conventions ensures that code intent is clear and reduces ambiguity during further development and review.

\subsection{Writing comments and Javadoc}

Documentation of the public interfaces was handled using Javadoc comments, following the
recommendations in Chapter~6 of the course book, in order to generate clear Java API documentation. Each public class and method includes a Javadoc description explaining its purpose, parameters, return values, and possible exceptions at a high level, without exposing internal implementation details.

Inline comments inside methods were deliberately avoided, since they tend to focus on internal implementation details rather than overall intent. Instead, clarity was achieved by using well-named methods and variables, allowing implementation details to remain encapsulated and easily modifiable without changing public interfaces.

\subsection{Code smell and refactoring}

During development, code smells such as duplicated code, long methods, and excessive use of primitive values were identified. When such issues were identified, the code was refactored to improve readability, cohesion, and maintainability without changing program behavior.

Refactoring was performed incrementally after implementing each system operation. This
continuous improvement process helped ensure that the code remained clean and aligned with
object-oriented design principles as the program evolved.

One common code smell is duplicated code, where similar logic appears in different methods. This can be resolved by extracting the shared functionality logic into separate helper methods, some are private methods if these are not used elsewhere. For instance, the formatting styles are common between the class \texttt{RepairOrder} and \texttt{RepairOrderDTO}, whereas the \texttt{format} function, that is only responsible for formatting the repair order string, can be extracted and placed in \texttt{RepairOrderDTO} class.

Another common code smell is long methods, which can be defined as a method taking responsibility of more than one functionality without explicit hierarchy of method calling. This can be resolved by splitting up each logic of the method that concerns each repsonsibility into its intended well-named method. This improves readability and makes the control flow easier to follow and understand. 

Additionally, primitive obsession is another code smell, where primitive values are favoured over objects. This can be adressed by introducing small data classes, DTOs, that wrap up multiple variables into one, such as \texttt{Cost} class that holds both an amount value and a currency value. This avoids primitive variables and fields to be scattered across the logic, which does not only cause low readability, but also harder to maintain since multiple variables have to be considered without an explicit relationship written as code syntax.


\subsection{Testing Strategy}

Unit testing was integrated into the development process and performed iteratively throughout implementation. After completing each system operation, relevant unit tests were written and executed before proceeding further.

Test cases were selected to cover typical usage scenarios as well as boundary cases, such as successful and unsuccessful lookup operations. Tests were implemented using the JUnit framework and focused on verifying the functional correctness of the controller, model, and integration layers. Methods that only produced output to \texttt{System.out} were excluded from testing in accordance with the seminar requirements.

This iterative testing approach helped detect errors early and ensured that newly added functionality did not break existing behavior.


% To be removed !!!
\iffalse
## 6.4 Code Smell and Refactoring
**Code smells** signal poor design and should be removed through refactoring to improve code quality without changing behavior.

* **Duplicated Code**: Duplicated code is a severe code smell that must be eliminated because even a single repeated statement can lead to errors, inefficiency, and maintenance disasters.
* **Long Method**:
  * A method is too long if its name cannot fully explain it — use **smaller, well‑named** methods instead of comments to make code clearer.
  * Reducing method calls does **not meaningfully improve performance**, so optimizing by avoiding them is pointless.
* **Large Class**: A class is too large when its cohesion is low — splitting related fields into a new class improves clarity and design.
* **Long Parameter List**:
  * Long parameter lists are a code smell — **replace primitive parameters with objects** to improve cohesion, encapsulation, and maintainable design. Sometimes the object itself has already contains the paramter, use **this** to get it.
  * Remove unnecessary parameters by letting the method obtain the needed data itself instead of passing it in.
* **Excessive Use of Primitive Variables**: Primitive Obsession is a code smell—use objects instead of excessive primitive fields to improve cohesion, encapsulation, and overall design quality.
  * **Group** closely related fields and methods into a new class to improve cohesion and ensure data is always valid.
  * Never use arrays to store unrelated primitive values—replace them with objects so each piece of data has a clear, meaningful name.
  * Use **enumerations** instead of strings or integers to avoid errors and make code meanings explicit.
* **Meaningless Names**: All declared program elements must have meaningful names. (packages, classes, interfaces, methods, fields, parameters, local variables, etc.).
  * **Avoid one‑letter identifiers**
      *   Prefer `Person person = new Person();` over `Person p = new Person();`.
      *   **Exceptions:**
          *   When the abstraction is naturally one letter (e.g., coordinate `x`).
          *   Loop variables (`i`, `j`, `k`, …) in nested loops.
  *   **Do not use meaningless temporary names like `tmp` or `temp`**
        *   Use descriptive names even for short‑lived variables (e.g., during swaps).
  *   **Never distinguish identifiers by adding numbers**
        *   Use meaningful names such as `fromAccount` and `toAccount`, **not** `account1` and `account2`.
  *   **Do not fear long names**
        *   What matters is clarity.
        *   Example: `firstCustomerBuyingCampaignItem` is acceptable if it clearly explains the variable’s purpose.
* **Anonymous Values**:
  *   **Always give every value an explaining name**
      *   Unnamed values (literals) make code hard to understand and maintain.
      *   Named values are safer than comments because the compiler enforces correctness, while comments can become outdated.
  *   **Do not place raw numeric values directly in method calls**
      *   Example: avoid `connect(10000)`.
      *   Prefer meaningful named constants like:
          *   `private static final int MILLIS_PER_SECOND = 1000;`
          *   `private static final int CONNECT_TIMEOUT_SECS = 10;`
      *   Use them in expressions such as:
          *   `connect(CONNECT_TIMEOUT_SECS * MILLIS_PER_SECOND);`
  *   **This naming practice applies to all primitive values—including strings**
      *   Example: avoid raw string values like `"\""` when building file paths.
      *   Use a constant:
          *   `private static final String PATH_SEPARATOR = "\\";`
      *   Then write:
          *   `openFile(dirName + PATH_SEPARATOR + fileName);`
  *   **Sometimes a method provides the clearest name for a value**
      *   Especially useful in conditional expressions (`if` statements).
      *   Example: instead of checking `if (c == 10)` for Unix end‑of‑line, create:
          *   `boolean isUnixEol(char c)`
      *   Use it as:
          *   `if (isUnixEol(ch))`
* **Complicated Flow Control**: Poorly structured **if‑statements and loops** become confusing when deeply nested or filled with unnecessary flags, making the logic hard to understand.
\fi


\section{Result}
\label{sec:result}

% \textbf{This chapter explains \textit{the result} of what you did.} \textbf{The report must show that you have done the work yourself and that you have understood what you have done}, both of these goals are met by explaining your solution here in the result chapter. Include links to your code in your Git repository (for seminars 3 and 4), and also include diagrams, see Figure \ref{fig:diag}, and other figures to illustrate your reasoning. All figures must be referenced in the text. Ask yourself if the solution is clearly explained, and if the reader will understand the application. What would you yourself want to know if you read about the application, is that included in the report? More specific instructions for the content can be found in each assignment.

Based on the design diagrams from seminar 2, the code implementation covers the basic flow and startup sequence. Additionally, the code implementation is structurally divided into multiple packages, where each of them corresponds to a layer in the MVC-layered architecture. 

Each system operation is described in terms of its purpose and execution flow through the different layers of the system. Relevant parts of the implementation are illustrated with code snippets, and the interaction between the controller, model, and integration layers is explained.

The complete source code for this seminar is available in the public Git repository at: \texttt{https://github.com/mo-mosverkstad/java-learning/tree/main/maven-projects/ElectricBikeRepair/}.

\subsection{\texttt{startup} sequence}

The startup sequence involves instantiation of the necessary entity class objects of the program, which can be mapped to the communication diagram of the start sequence. The sequence is important during development as the first to implement, since otherwise the program does not have a clear starting point and can be harder to track as the code grows large. Additionally, it also provides a possibility for manually checking program behavior during implementation. 

The startup sequence initializes the View, Controller, and Registry classes, which together form the core structure of the application. At the top of the layered MVC architecture, the View class provides a simplified interface with hard-coded input and output to the command line. Since there are two different roles, two separate classes \texttt{ReceptionistController} and \texttt{TechnicianController} are created, according to the design. However, some of their authorization actions overlap, both classes share a common superclass \texttt{Controller}.
Both controllers are initialized and handled collectively by a \texttt{ControllerCreator} class. In this case, only the creator class needs to be initialized instead of separately initializing each class. As number of roles as well as different controllers grows, bundling together is easier to extend and maintain than keeping them as separate classes.
In integration layer, the registry classes handle the data registry. Since the code does not handle databases, the implementation will store the data in-memory instead.
The registry classes consist of \texttt{CustomerRegistry} and \texttt{RepairOrderRegistry}. \texttt{CustomerRegistry} manages customer data and detailed information, while \texttt{RepairOrderRegistry} manages repair orders as well as their details. Both registry classes are collectively handled by a \texttt{RegistryCreator}, which is a class that manages both registry classes, with the same pattern as for handling controllers.

In this case, the fundamental functional skeleton of the program can be placed out, whereas each system operation can easily be added by extending the existing structure.

This snippet shows how the application is started and how the different layers are connected. The startup logic is centralized and contains no business logic, which keeps object creation separated from domain behavior.

The main class initializes all the class object and acts as the application entry point. It establishes the overall program structure by creating instances of the integration, controller, and view layers, while connecting them together as a starting point integration.

\begin{lstlisting}[style=javastyle, caption={Main class and main method entry point}]
public class Main {
    public static void main(String[] args) {
        RegistryCreator registryCreator = new RegistryCreator();
        Printer printer = new Printer();
        ControllerCreator controllerCreator = new ControllerCreator(registryCreator, printer);
        View view = new View(controllerCreator);
        view.proceedActions("0707654321", "SK-2024-055");
    }
}
\end{lstlisting}

The View class currently contains a dummy view interface with hardcoded input as well as output to command line interface by using system print.

Actions proceeded by receptionist and technician are hardcoded and collected in the method \texttt{proceedActions} that calls the internal private functions. Internal functions \texttt{proceedReceptionPreparationActions} and \texttt{proceedReceptionConfirmationActions} are proceeded by receptionist using \texttt{ReceptionistController} while \texttt{proceedTechnicianDiagnosticActions} proceeded by technician with \texttt{TechnicianController}-

\begin{lstlisting}[style=javastyle, caption={Skeleton code of \texttt{View} class}]
package se.ebikerepair.view;

/**
 * View layer that orchestrates the repair workflow: reception preparation, technician diagnostics,
 * and reception confirmation. Catches and displays user-facing errors.
 */
public class View {
    private final ReceptionistController receptionistController;
    private final TechnicianController technicianController;

    /**
     * Creates a view with controllers from the given controller creator.
     *
     * @param controllerCreator the factory providing access to controllers
     */
    public View(ControllerCreator controllerCreator) {
        receptionistController = controllerCreator.getReceptionistController();
        technicianController = controllerCreator.getTechnicianController();
    }

    /**
     * Runs the full repair workflow: reception preparation, technician diagnostics, and reception confirmation.
     *
     * @param telephoneNumber the customer's telephone number
     * @param bikeSerialNumber the serial number of the broken bike
     */
    public void proceedActions(String telephoneNumber, String bikeSerialNumber) {
        proceedReceptionPreparationActions(telephoneNumber, bikeSerialNumber);
        proceedTechnicianDiagnosticActions(telephoneNumber);
        proceedReceptionConfirmationActions(telephoneNumber);
    }
    ...
}
\end{lstlisting}

\subsection{The \texttt{searchCustomer} system operation}

The first implemented system operation is \texttt{searchCustomer}, which involves retrieving customer information based on a provided telephone number. This operation is initiated from the View layer and handled by the \texttt{ReceptionistController}, following the interaction defined in the communication diagram.

The execution of this operation begins when the View calls \texttt{searchCustomer(telephoneNumber)} in the controller. The controller normalizes the input telephone number into E.164 format using the \texttt{TelephoneNumber} utility class. The number is used as a key to find the customer in \texttt{CustomerRegistry}, which returns the corresponding \texttt{CustomerDTO}, or \texttt{null} if no match is found. In the end, the controller validates the result: If a customer is found, it is returned to the caller. Otherwise, an \texttt{IllegalArgumentException} is thrown.

This applies the separation of concerns, where the input normalization is handled by a utility class, outside of the main logic. Additionally, data access is handled by the registry, while the control flow and validation are handled by the controller. In terms of encapsulation, data structures of registry are hidden from controller. 

\begin{lstlisting}[style=javastyle, caption={Method \texttt{proceedReceptionPreparationActions} in \texttt{View} for proceeding receptionist preparation tasks}]
private void proceedReceptionPreparationActions(String telephoneNumber, String bikeSerialNumber) {
    try {
        CustomerDTO foundCustomer = receptionistController.searchCustomer(telephoneNumber);
        printout("1. Reception - Found customer:", foundCustomer);
    } catch (IllegalArgumentException | IllegalStateException e) {
        System.out.println(ERROR_PREFIX + e.getMessage());
    }
}
\end{lstlisting}

This code illustrates how input normalization, data access, and validation responsibilities are clearly separated. The controller coordinates the operation, while telephone formatting and data storage are delegated to specialized classes.

\begin{lstlisting}[style=javastyle, caption={Searching a customer using \texttt{searchCustomer} method in \texttt{ReceptionistController}}]
public class ReceptionistController extends Controller{
    private final CustomerRegistry customerRegistry;
    private final Printer printer;

    /**
     * Searches for a customer by telephone number.
     *
     * @param telephoneNumber the customer's telephone number (any Swedish format)
     * @return the found customer DTO
     * @throws IllegalArgumentException if the phone number format is invalid or customer not found
     */
    public CustomerDTO searchCustomer(String telephoneNumber) throws IllegalArgumentException{
        String phoneNumberE164 = new TelephoneNumber(telephoneNumber).toE164();
        CustomerDTO foundCustomer = customerRegistry.find(phoneNumberE164);
        if (foundCustomer == null) {
            throw new IllegalArgumentException("No customer found for telephone number: " + telephoneNumber);
        }
        return foundCustomer;
    }
}
\end{lstlisting}

\subsection{The \texttt{createRepairOrder} system operation}

The \texttt{createRepairOrder} operation is responsible for creating a new repair order for a specific customer and their bike. The operation is initiated by the View after a customer has been successfully retrieved. The View provides the telephone number and a \texttt{ProblemDTO}, which contains both the bike and a description of the issue.

To start with, the View calls \texttt{createRepairOrder} in the \texttt{ReceptionistController}. The controller retrieves the customer using the telephone number, as a new \texttt{RepairOrder} object is created in the model layer. The problem description is added to the repair order, which is stored in the \texttt{RepairOrderRegistry}. In the end, the unique repair order id is returned back to the View

This operation shows how responsibilities are distributed across layers. The controller coordinates the process while the model encapsulates business logic, and the registry manages persistence. By generating a unique id for each repair order, the system ensures traceability and enables future operations such as diagnostics and task management.

As a result, a new repair order is successfully created and stored, allowing the workflow to proceed to the diagnostic phase.

\begin{lstlisting}[style=javastyle, caption={Method \texttt{proceedReceptionPreparationActions} in \texttt{View} for proceeding receptionist preparation tasks}]
private void proceedReceptionPreparationActions(String telephoneNumber, String bikeSerialNumber) {
    try {
        CustomerDTO foundCustomer = receptionistController.searchCustomer(telephoneNumber);
        printout("1. Reception - Found customer:", foundCustomer);

        BikeDTO foundBike = foundCustomer.getBikeBySerialNumber(bikeSerialNumber);
        String repairOrderId = receptionistController.createRepairOrder(telephoneNumber, new ProblemDTO("Broken bike chain", foundBike));
        printout("2. Reception - Created repair order with id: ", repairOrderId);

    } catch (IllegalArgumentException | IllegalStateException e) {
        System.out.println(ERROR_PREFIX + e.getMessage());
    }
}
\end{lstlisting}

The detailed implementation of \texttt{createRepairOrder} starts by finding a customer by a given telephone number, while a \texttt{RepairOrder} instance is created. The repair order is then saved into the \texttt{repairOrderRegistry}.

\begin{lstlisting}[style=javastyle, caption={Creating a new repair order using \texttt{createRepairOrder} method in \texttt{ReceptionistController}}]
/**
 * Creates a new repair order for the customer identified by telephone number.
 *
 * @param telephoneNumber the customer's telephone number (any Swedish format)
 * @param problemDTO the problem description with the broken bike
 * @return the generated repair order ID
 * @throws IllegalArgumentException if the customer is not found or phone format is invalid
 */
public String createRepairOrder(String telephoneNumber, ProblemDTO problemDTO) throws IllegalArgumentException{
    CustomerDTO foundCustomer = searchCustomer(telephoneNumber);
    RepairOrder repairOrder = new RepairOrder(foundCustomer);
    repairOrder.updateProblem(problemDTO);
    repairOrderRegistry.save(repairOrder);
    return repairOrder.getId();
}
\end{lstlisting}

\subsection{The \texttt{findRepairOrder} system operation}

The \texttt{findRepairOrder} operation is responsible for retrieving an existing repair order associated with a specific customer. This operation is essential for continuing the workflow after a repair order has been created, as it allows subsequent actions such as diagnostics and task creation to be performed on the correct entity.

The operation is initiated from the View layer by providing a telephone number and is handled by the controller. The controller first normalizes the telephone number into E.164 format to ensure consistency. It then retrieves all repair order identifiers associated with the customer from the \texttt{RepairOrderRegistry}. From this set, the most recent repair order is selected and fetched from the registry. The retrieved repair order is converted into a \texttt{RepairOrderDTO} before being returned to the View.

If no repair orders exist for the given customer, an \texttt{IllegalArgumentException} is thrown. In the case of data inconsistency, where an identifier exists without a corresponding repair order, an \texttt{IllegalStateException} is raised.

This operation demonstrates how the system maintains continuity in the workflow by consistently retrieving the most recent repair order. The use of DTOs ensures that internal model data is not exposed directly to the View, preserving encapsulation. As a result, the system provides access to the current state of a repair order, enabling further processing in subsequent operations.

\begin{lstlisting}[style=javastyle, caption={Method \texttt{proceedTechnicianDiagnosticActions} in \texttt{View} for proceeding receptionist preparation tasks}]
private void proceedTechnicianDiagnosticActions(String telephoneNumber) {
    try {
        RepairOrderDTO repairOrderDTO = technicianController.findRepairOrder(telephoneNumber);
        printout("3. Technician - Requested repair order: ", repairOrderDTO);
    } catch (IllegalArgumentException | IllegalStateException e) {
        System.out.println(ERROR_PREFIX + e.getMessage());
    }
}
\end{lstlisting}

The \texttt{findRepairOrder} operation retrieves the most recently created repair order associated with a customer.

\begin{lstlisting}[style=javastyle, caption={Method for finding a repair order by a given telephone number \texttt{findRepairOrder} in \texttt{Controller} (shared by both technician and receptionist)}]
/**
 * Finds the most recent repair order for the given telephone number.
 *
 * @param telephoneNumber the customer's telephone number (any Swedish format)
 * @return the most recent repair order as a DTO
 * @throws IllegalArgumentException if no repair orders exist for the customer or phone format is invalid
 * @throws IllegalStateException if the repair order ID exists but the order cannot be found (data inconsistency)
 */
public RepairOrderDTO findRepairOrder(String telephoneNumber) throws IllegalArgumentException, IllegalStateException{
    List<String> repairOrderIds = findRepairOrderIds(telephoneNumber);
    if (repairOrderIds.isEmpty()) {
        throw new IllegalArgumentException("There is no repair order created for this customer.");
    }
    String id = repairOrderIds.get(0);
    RepairOrder repairOrder = repairOrderRegistry.findByRepairOrderId(id);
    if (repairOrder == null) {
        throw new IllegalStateException("Repair order not found for id: " + id);
    }
    return repairOrder.toDTO();
}
\end{lstlisting}

\subsection{The \texttt{addDiagnosticResult} system operation}

The \texttt{addDiagnosticResult} operation allows a technician to record the outcome of a diagnostic task for a specific repair order. This operation represents the inspection phase of the workflow, where the condition of the bike is evaluated and documented.

The operation is initiated from the View layer after a repair order has been retrieved. The technician provides the repair order identifier, the name of the diagnostic task, and a \texttt{ResultDTO} describing the outcome. The controller retrieves the corresponding \texttt{RepairOrder} from the \texttt{RepairOrderRegistry} and delegates the update to the model layer. The \texttt{RepairOrder} then updates its internal \texttt{DiagnosticReport} by locating the specified task and applying the result. Once the update is complete, the modified repair order is saved back into the registry.

The responsibility for modifying diagnostic data is encapsulated within the model, ensuring that business rules remain centralized and protected from external layers. The controller acts only as a coordinator of the operation, without directly manipulating internal data structures.

As a result, the repair order is updated with new diagnostic information, which directly influences later stages of the workflow such as repair task creation and cost estimation.

\begin{lstlisting}[style=javastyle, caption={Method \texttt{proceedTechnicianDiagnosticActions} in \texttt{View} for proceeding receptionist preparation tasks}]
private void proceedTechnicianDiagnosticActions(String telephoneNumber) {
    try {
        RepairOrderDTO repairOrderDTO = technicianController.findRepairOrder(telephoneNumber);
        printout("3. Technician - Requested repair order: ", repairOrderDTO);

        String repairOrderId = repairOrderDTO.id();
        String diagnosticTaskName = "Mechanical Safety Check";
        ResultDTO result = new ResultDTO(true, true, "Chain and gears should be replaced.");
        technicianController.addDiagnosticResult(repairOrderId, diagnosticTaskName, result);
    } catch (IllegalArgumentException | IllegalStateException e) {
        System.out.println(ERROR_PREFIX + e.getMessage());
    }
}
\end{lstlisting}

The detailed \texttt{addDiagnosticResult} is implemented by adding a diagnostic result to a tied diagnostic task, which is identified using a repair order id.

\begin{lstlisting}[style=javastyle, caption={Add diagnostic result using \texttt{addDiagnosticResult} method in \texttt{TechnicianController}}]
/**
 * Adds a diagnostic result to a diagnostic task identified by name within a repair order.
 *
 * @param repairOrderId the ID of the repair order
 * @param diagnosticTaskName the name (partial match) of the diagnostic task to update
 * @param result the result to apply to the diagnostic task
 * @throws IllegalStateException if the repair order is not found
 * @throws IllegalArgumentException if no diagnostic task matches the given name
 */
public void addDiagnosticResult(String repairOrderId, String diagnosticTaskName, ResultDTO result) throws IllegalStateException, IllegalArgumentException {
    RepairOrder repairOrder = repairOrderRegistry.findByRepairOrderId(repairOrderId);
    if (repairOrder == null) {
        throw new IllegalStateException("Repair order not found for id: " + repairOrderId);
    }
    repairOrder.updateDiagnosticResult(diagnosticTaskName, result);
    repairOrderRegistry.save(repairOrder);
}
\end{lstlisting}

\subsection{The \texttt{createProposedRepairTask} system operation}

The \texttt{createProposedRepairTask} operation allows a technician to define repair actions based on the results of the diagnostic phase. This operation represents the transition from inspection to repair planning, where identified issues are translated into concrete repair tasks.

The operation is initiated from the View layer after diagnostic results have been recorded. The technician provides a \texttt{RepairTaskDTO} containing a description of the repair, along with its associated cost and estimated duration. The controller retrieves the corresponding \texttt{RepairOrder} from the registry and delegates the responsibility of adding the task to the model layer. The \texttt{RepairOrder} then adds the new task to its \texttt{RepairTaskCollection}, ensuring that it becomes part of the overall repair plan. The updated repair order is subsequently stored in the registry.

The model layer manages the internal collection of repair tasks, ensuring consistency and preventing unauthorized access to internal data structures. This encapsulation allows the system to evolve without affecting higher-level components.

As a result, the repair order is extended with proposed tasks, forming the basis for customer approval and enabling the next stage of the workflow.

\begin{lstlisting}[style=javastyle, caption={Internal method \texttt{proceedTechnicianDiagnosticActions} in \texttt{View} for proceeding receptionist preparation tasks}]
private void proceedTechnicianDiagnosticActions(String telephoneNumber) {
    try {
        RepairOrderDTO repairOrderDTO = technicianController.findRepairOrder(telephoneNumber);
        printout("3. Technician - Requested repair order: ", repairOrderDTO);

        String repairOrderId = repairOrderDTO.id();
        String diagnosticTaskName = "Mechanical Safety Check";
        ResultDTO result = new ResultDTO(true, true, "Chain and gears should be replaced.");
        technicianController.addDiagnosticResult(repairOrderId, diagnosticTaskName, result);
        printout("4. Technician - Updated diagnostic task: ", diagnosticTaskName);

        RepairTaskDTO repairTask1 = new RepairTaskDTO("Replace Chain", "Removal of worn or stretched chain and installation of a new compatible e\u2011bike chain. Includes lubrication, tension adjustment, and drivetrain alignment check.", new Cost(500, "SEK"), 1);
        RepairTaskDTO repairTask2 = new RepairTaskDTO("Replace gears", "Replacement of worn or damaged rear cassette/freewheel or front chainrings. Includes removal of old components, installation of new gear set, indexing and tuning of derailleur(s).", new Cost(400, "SEK"), 2);
        technicianController.addRepairTask(repairOrderId, repairTask1);
        technicianController.addRepairTask(repairOrderId, repairTask2);
        printout("5. Technician - Created proposed repair tasks 01: ", repairTask1);
        printout("6. Technician - Created proposed repair tasks 02: ", repairTask2);
    } catch (IllegalArgumentException | IllegalStateException e) {
        System.out.println(ERROR_PREFIX + e.getMessage());
    }
}
\end{lstlisting}

Detailed implementation of \texttt{addRepairTask} method attaches a specific repair task onto a repair order.

\begin{lstlisting}[style=javastyle, caption={Add repair task using \texttt{addRepairTask} method in \texttt{TechnicianController}}]
    /**
     * Adds a repair task to a repair order.
     *
     * @param repairOrderId the ID of the repair order
     * @param repairTask the repair task to add
     * @throws IllegalStateException if the repair order is not found
     */
    public void addRepairTask(String repairOrderId, RepairTaskDTO repairTask) throws IllegalStateException {
        RepairOrder repairOrder = repairOrderRegistry.findByRepairOrderId(repairOrderId);
        if (repairOrder == null) {
            throw new IllegalStateException("Repair order not found for id: " + repairOrderId);
        }

        repairOrder.addRepairTask(repairTask);
        repairOrderRegistry.save(repairOrder);
    }
\end{lstlisting}

\subsection{The \texttt{acceptRepairOrder} and \texttt{rejectRepairOrder} system operation}

The \texttt{acceptRepairOrder} operation is performed when a customer accepts the proposed repair tasks.

The operation is initiated from the View layer after the proposed repair tasks have been presented to the customer. The controller retrieves the corresponding \texttt{RepairOrder} from the registry and updates its state to \texttt{Accepted}. This state transition is handled within the model, ensuring that the lifecycle of the repair order remains consistent. After updating the state, the controller sends the repair order to the \texttt{Printer}, producing a formal output. The updated repair order is then saved in the registry.

As a result, the repair order is marked as accepted and confirmed, ready for execution while a formal record of the accepted repair is generated.

\begin{lstlisting}[style=javastyle, caption={Internal method \texttt{proceedReceptionConfirmationActions} in \texttt{View} that handles receptionist confirmation tasks}]
private void proceedReceptionConfirmationActions(String telephoneNumber) {
    try {
        RepairOrderDTO repairOrderDTO = receptionistController.findRepairOrder(telephoneNumber);
        printout("7. Reception - found repair order: ", repairOrderDTO);

        printout("8. Reception - Accepted order");
        String repairOrderId = repairOrderDTO.id();
        receptionistController.acceptRepairOrder(repairOrderId);

    } catch (IllegalArgumentException | IllegalStateException e) {
        System.out.println(ERROR_PREFIX + e.getMessage());
    }
}
\end{lstlisting}

Internal implementation of \texttt{acceptRepairOrder} marks the repair order as accepted.

\begin{lstlisting}[style=javastyle, caption={Implementation of \texttt{acceptRepairOrder} in \texttt{ReceptionistController}}]
/**
 * Accepts a repair order, updates its state, saves it, and prints it.
 *
 * @param repairOrderId the ID of the repair order to accept
 * @throws IllegalStateException if the repair order is not found
 */
public void acceptRepairOrder(String repairOrderId) throws IllegalStateException{
    RepairOrder repairOrder = repairOrderRegistry.findByRepairOrderId(repairOrderId);
    if (repairOrder == null) {
        throw new IllegalStateException("Repair order not found for id: " + repairOrderId);
    }
    repairOrder.accept();
    repairOrderRegistry.save(repairOrder);
    printer.print(repairOrder);
}
\end{lstlisting}

The \texttt{rejectRepairOrder} operation is used when a customer declines the proposed repair tasks. This operation represents an alternative outcome in the workflow, where the repair process is not approved.

The operation is initiated from the View layer when the customer decides not to proceed. The controller retrieves the corresponding \texttt{RepairOrder} from the registry and updates its state to \texttt{Rejected}. This state change is performed within the model layer, ensuring that all lifecycle transitions are handled consistently. The updated repair order is then saved back into the registry.

By supporting both acceptance and rejection scenarios, the system provides a complete representation of the decision-making process within the repair workflow. The encapsulation of state changes within the model ensures that the system remains robust and maintainable.

As a result, the repair order is marked as rejected, and no further actions are performed on it.

\begin{lstlisting}[style=javastyle, caption={Implementation of \texttt{rejectRepairOrder} in \texttt{ReceptionistController}}]
/**
 * Rejects a repair order and updates its state.
 *
 * @param repairOrderId the ID of the repair order to reject
 * @throws IllegalStateException if the repair order is not found
 */
public void rejectRepairOrder(String repairOrderId) throws IllegalStateException{
    RepairOrder repairOrder = repairOrderRegistry.findByRepairOrderId(repairOrderId);
    if (repairOrder == null) {
        throw new IllegalStateException("Repair order not found for id: " + repairOrderId);
    }
    repairOrder.reject();
    repairOrderRegistry.save(repairOrder);
}
\end{lstlisting}

\subsection{Result printouts}
In this chapter, the JUnit test result and program execution result are illustrated.

\subsubsection{Testing printout}
The testing command \texttt{mvn compile test} produces the testing summary \ref{fig:sample-test} of tests.

\begin{lstlisting}[style=outputstyle, caption={JUnit test execution summary}, label={fig:sample-test}]
...
[INFO]
[INFO] Results:
[INFO]
[INFO] Tests run: 116, Failures: 0, Errors: 0, Skipped: 0
[INFO]
...
\end{lstlisting}

\subsubsection{Sample program execution printout}
Figure~\ref{fig:sample-run} shows a sample execution of the program, demonstrating the basic workflow from customer lookup to repair order acceptance. The output is produced by executing the program using \texttt{mvn compile exec:java -Dexec.mainClass="se.ebikerepair.startup.Main"}.

\begin{lstlisting}[style=outputstyle, caption={Sample output from program execution}, label={fig:sample-run}]
1. Reception - Found customer:
=================================
  Customer Information
=================================
  Name:      Astrid Johansson
  Phone:     +46707654321
  Email:     astrid.johansson@example.com
  Bikes:
    - Brand: Monark, Model: E-Karin (S/N: MO-2024-010)
    - Brand: Skeppshult, Model: Elit (S/N: SK-2024-055)
=================================

2. Reception - Created repair order with id: 
f6ff76e9-bc5c-4a68-96f8-5e3223194795

3. Technician - Requested repair order: 
...

8. Reception - Accepted order

**** Printing repair order ****
=================================
  Repair Order
=================================
  Order ID:       f6ff76e9-bc5c-4a68-96f8-5e3223194795
  Status:         Accepted
  Created:        Sat Apr 18 14:51:24 CEST 2026
  Est. Complete:  Wed Apr 22 14:51:24 CEST 2026
  Total Cost:     1175.00 SEK
********************************
  Customer:
  - Name:      Astrid Johansson
  - Phone:     +46707654321
  - Email:     astrid.johansson@example.com
  - Bikes:
    - Brand: Monark, Model: E-Karin (S/N: MO-2024-010)
    - Brand: Skeppshult, Model: Elit (S/N: SK-2024-055)
********************************
  Problem:
    - Description: Broken bike chain
    - Broken bike: Skeppshult Elit (S/N: SK-2024-055)
********************************
  Diagnostic Report:
    Days: 1
    Cost: 275.00 SEK
    Description: This report contains pre-defined diagnostic tasks for e-bike inspection. Only tasks marked with [X] have been performed and contribute to the repair cost.
      --------------------------------
      [ ] 0.Basic Electrical System Check
            A general inspection of the e-bike's electrical components,
            including wiring, connectors, sensors, and display functionality.
            Checks for loose connections, corrosion, and obvious cable damage.
      [ ] Cost: 300.00 SEK | Est. days: 1 | Result: have not checked yet
      --------------------------------
      [ ] 1.Battery Health & Capacity Test
            Tests the battery's remaining capacity, charge cycles, internal
            resistance, and voltage output under load. Identifies early signs of
            cell degradation.
      [ ] Cost: 425.00 SEK | Est. days: 2 | Result: have not checked yet
      --------------------------------
      [ ] 2.Motor Diagnostic & Noise Analysis
            Evaluates motor performance, noise levels, internal resistance,
            torque output, error codes, and abnormal vibration. Includes hub
            motors or mid-drive systems.
      [ ] Cost: 500.00 SEK | Est. days: 2 | Result: have not checked yet
      --------------------------------
      [ ] 3.Charger Functionality Test
            Measures charger voltage, consistency, output under load, and
            connector integrity to ensure safe and stable charging.
      [ ] Cost: 200.00 SEK | Est. days: 1 | Result: have not checked yet
      --------------------------------
      [X] 4.Drivetrain & Mechanical Safety Check
            Ensures gears, chain, brakes, bearings, and frame components work
            correctly with the electric assist system. Identifies mechanical
            issues affecting motor load.
      [TO BE REPAIRED] Cost: 275.00 SEK | Est. days: 1 | Result: Chain and gears should be replaced.
********************************
  Proposed Repair Tasks:
    Days: 3
    Cost: 900.00 SEK
    --------------------------------
    [X] Replace Chain
        Removal of worn or stretched chain and installation of a new compatible
        e-bike chain. Includes lubrication, tension adjustment, and drivetrain
        alignment check.
         Cost: 500.00 SEK | Est. days: 1
    --------------------------------
    [X] Replace gears
        Replacement of worn or damaged rear cassette/freewheel or front
        chainrings. Includes removal of old components, installation of new gear
        set, indexing and tuning of derailleur(s).
         Cost: 400.00 SEK | Est. days: 2
=================================
\end{lstlisting}

\section{Discussion}

% \textbf{This chapter evaluates your solution} according to the assessment criteria found in the assessment-criteria documents, which are available under each seminar (sections 4.1-4.4) on the \texttt{Seminar Tasks} page of the course web. To get 2p, it's important that the evaluation is extensive, and \textbf{motivated with concrete examples from your solution}. General statements, that could be applied to any solution, will not contribute to a higher score.

\section*{References}

\textbf{It's mandatory to list here all sources from where you have collected information}, for example course material and YouTube videos. \textbf{Remember that all AI usage must be specified}. There are examples of how AI usage can be specified in Canvas, on the course front page.

You can use any reference style you wish. \textbf{Don't mention other students you have collaborated with}, that shall instead be mentioned in the \textit{Project Members} section at the top of the report.

\end{document}
